\documentclass[11pt]{article}

\usepackage[margin=1in]{geometry}
\usepackage{amsmath}
\usepackage{booktabs}
\usepackage{enumitem}
\usepackage{graphicx}
\usepackage{siunitx}
\usepackage{hyperref}
\usepackage{gensymb}


\title{Basic Crawling Test}
\author{Rift / Team 09}
\date{\today}

\begin{document}


\maketitle

\section*{Test ID}
\textbf{T-01BC Basic Crawling Test}

\section*{Test Objective}
Quantify the traversal speed of the rover on flat ground, over a set number of steps. A node path will be planned for a set number of steps, where the distance and time used for that traversal will be measured.  Speed will be measured for both forward-backward and side-to-side traversal. Quantify the average node height for forward flat ground traversal.

\section*{Robot Configuration}
\begin{itemize}[noitemsep]
    \item Robot configurations: Standing
    \item Tube pressures: 15 psi
    \item Control mode: Open-loop
    \item Electronics status: Batteries are fully charged
\end{itemize}

\section*{Materials}
\begin{itemize}[noitemsep]
    \item Stopwatch timer
    \item Measuring tape (at least 5 m in length)
    
\end{itemize}

\section*{Environment and Setup}
\begin{itemize}[noitemsep]
    \item Surface type: Finished Smooth Concrete
    \item Surface angle: $\psi = 0 {\degree}$ 
    \item External load: None 
    \item Measurement method: measuring time to complete path and distance traveled; measuring average height of bottom node during traversal
\end{itemize}

\section*{Initial Conditions}
\begin{itemize}[noitemsep]
    \item Initial position: $x_0 = 0$
    \item Initial orientation: $\theta_z = 0 {\degree}$
    \item Initial shape state: Standing
    \item Stabilization time before command: 10?~s
\end{itemize}

\section*{Command Sequence}
\begin{itemize}[noitemsep]
    \item Command type: \underline{\hspace{5 cm}}
    \item Command values: \underline{\hspace{5 cm}}
    \item Update rate: 3~Hz
    \item Command Steps: 8~s
\end{itemize}

\section*{Measured Quantities}
\begin{itemize}[noitemsep]
    \item Forward traversal speed
    \item Side-to-side traversal ("crab walking")
    \item Average node height during forward traversal compared to predicted node height in simulation. 
\end{itemize}

\section*{Success Criteria (Per Trial)}
\begin{itemize}[noitemsep]
    \item Forward traversal speed: .1 km/h  $\leq$ .14 km/h $\leq$ No upper limit
    \item Side to Side traversal speed: ?? km/h  $\leq$ ?? km/h $\leq$ No upper limit
    \item  Step height: 10 cm  $\leq$ 25 cm $\leq$ No upper limit
\end{itemize}

\section*{Test Procedure}
\begin{enumerate}[noitemsep]
    \item Place robot in standing position
    \item Complete pretest checklist
    \item Send command sequence to drive rover 8 steps (2 full step cycles) forward while starting timer
    \item Stop timer upon completion of sequence
    \item Measure distance between front left and front right nodes before and after step sequence.
    \item Record average of the two above measured distances divided by traversal time
    \item Repeat steps 1-7, driving the rover side-to-side instead of forward
    \item Reset robot to standing position
    \item Send command sequence to drive rover halfway through step
    \item Measure height of stepping node
    \item Repeat steps 9 and 10 through 8 sets of steps
    \item Record average of step heights

\end{enumerate}

\section*{Number of Trials}
\begin{itemize}[noitemsep]
    \item Total trials: 1
    \item Time between trials: N/A
\end{itemize}

\section*{Results}
Summarize the observed performance metrics.

\begin{center}
\begin{tabular}{lc}
\toprule
Metric & Value\\
\midrule
Front Traversal Speed (km/h)&  \\
 Side Traversal Speed (km/h)&\\
 Average Step Height (cm)&\\
\end{tabular}
\end{center}

\section*{Observed Failure Modes}
Describe any observed failure behaviors.

\begin{itemize}[noitemsep]
    \item Truss buckles during stepping
    \item Truss fails to lift nodes from the ground
    \item Truss tips while lifting node from the ground
\end{itemize}

\section*{Conclusions}
Provide a concise interpretation of the results.

\textit{Example:} The robot achieved consistent forward motion on flat ground with an average speed of \underline{\hspace{1 cm}}~cm/s, and a step height of  \underline{\hspace{1 cm}}~cm. 

\section*{Notes and Limitations}
Document sources of variability or limitations.

\end{document}